\section{Design}
\label{sec:design}

\subsection{System Overview}

\namei places a programmable decision at the point where the VFS chooses the
visible object for a pathname or directory entry. Four invariants organize the
design: (1) policy runs only at lookup and directory enumeration, (2) policy
returns bounded path-view actions over existing lower objects, (3) the kernel
validates every action before ordinary VFS execution resumes, and (4) the lower
filesystem retains permissions, file operations, data, writes, page cache,
persistence, and consistency.

Figure~\ref{fig:overview} shows the resulting flow. An application in a
workload scope performs an ordinary lookup, open, stat, exec, or readdir. The
VFS reaches the \namei decision point before committing to the visible object
for the current component or directory entry. The eBPF policy returns a bounded
path-view action, and ordinary VFS and lower-filesystem logic resumes after
kernel validation.

\begin{figure}[t]
\centering
\scriptsize
\setlength{\tabcolsep}{3pt}
\begin{tabular}{c c c c c}
\fbox{\begin{tabular}{c}workload\\cgroup\end{tabular}} &
\begin{tabular}{c}$\longrightarrow$\\lookup/readdir\end{tabular} &
\fbox{\begin{tabular}{c}VFS\\name resolution\end{tabular}} &
\begin{tabular}{c}$\longrightarrow$\\event\end{tabular} &
\fbox{\begin{tabular}{c}eBPF\\policy\end{tabular}} \\
&&&& $\downarrow$ \\
\fbox{\begin{tabular}{c}lower FS owns\\dentries, inodes, data, writes\end{tabular}} &
\begin{tabular}{c}$\longleftarrow$\\ordinary VFS path\end{tabular} &
\fbox{\begin{tabular}{c}kernel validates\\bounded action\end{tabular}} &
\begin{tabular}{c}$\longleftarrow$\\pass/redirect/hide/select\end{tabular} &
\fbox{\begin{tabular}{c}policy output\\no FS methods\end{tabular}}
\end{tabular}
\caption{\namei keeps policy at VFS name resolution. The policy chooses bounded
path-view actions, while the kernel and lower filesystem continue to own
filesystem objects and data semantics.}
\label{fig:overview}
\end{figure}

\subsection{Name-Resolution Boundary Goals}

Like \code{sched_ext}, \namei lets BPF choose policy while the kernel keeps the
VFS mechanisms that own objects and execute filesystem semantics. The VFS and
lower filesystem retain path walking, dentries, inodes, permissions, file
operations, page cache, writes, persistence, and consistency.

The design defines a small kernel extension for name-resolution policy while
leaving full filesystem expressiveness to FUSE, stackable, or custom
filesystems. Table~\ref{tab:req-design} maps the four design goals to the RQs.

The current action set allows pass-through, replacement-component redirect,
hide, and selection from kernel-registered lower targets. Access-control denial
remains outside the current action set because \namei focuses on selecting or
hiding existing objects during name resolution.

\begin{table}[t]
\centering
\scriptsize
\begin{tabular}{@{}p{0.19\linewidth}p{0.5\linewidth}p{0.17\linewidth}@{}}
\hline
Goal & Design choice & RQ \\
\hline
G1 Expressiveness & invoke one workload-scoped decision contract at lookup and directory-enumeration time & RQ1 \\
G2 Lower-filesystem ownership & return only bounded actions over existing lower objects while keeping data, writes, permissions, page cache, and persistence below the hook & RQ1, RQ3 \\
G3 Placement cost & run bounded policy at lookup/readdir instead of routing the same decision through a filesystem service & RQ2 \\
G4 Fail-closed safety & use verified eBPF outputs plus kernel validation so malformed or unsupported decisions fail visibly & RQ3 \\
\hline
\end{tabular}
\caption{Design-goal mapping. If an oracle needs synthetic contents, custom
metadata, or data-path mediation, the effect belongs outside \namei's
boundary.}
\label{tab:req-design}
\end{table}

\subsection{Policy Model}

The policy model follows directly from the workload-derived requirements
because policy must (1) see lookup and directory-enumeration events, (2) carry
workload state, (3) return only bounded path-view actions, and (4) let the
kernel validate the result.

Here, the transition is the workload state change, the policy is the eBPF
decision logic, and the action is the bounded result validated by the kernel.

The policy input contains operation context, the current component, parent
context, and policy state. The policy state is workload-defined but not a
filesystem object model. For the representative workloads, policy state is a
branch or checkpoint identifier in an agent workspace, a cache or validation
epoch in an environment run, or a configuration/secret epoch in a service.

The output passes the lower lookup, redirects selection to a replacement
component, hides a name from lookup and directory enumeration, or selects a
registered target. Directory enumeration uses the same decision function so
that aliases returned by readdir match the objects selected by lookup.

The kernel checks the decision, preserves normal permission and
lower-filesystem checks for the selected object, maintains ordinary VFS path
resolution invariants, and fails invalid decisions visibly instead of falling
back to a weaker policy.

\subsection{Filesystem Ownership Boundary}

\namei excludes filesystem ownership by leaving inode operations, file
operations, dentry and inode allocation, synthetic contents, custom metadata,
read/write mediation, distributed indexes, and cross-path transactions outside
the hook.

The name-resolution boundary helps only when the source transition can
be expressed as lookup selection or directory visibility. When the oracle
requires data-path mediation, synthetic content, custom metadata, or
application-level orchestration, the transition remains motivation or related
work rather than \namei evidence.
